\documentclass[UTF8,a4paper,11pt]{ctexart}

\usepackage{amsmath,amssymb}
\usepackage{booktabs}
\usepackage{geometry}
\usepackage{enumitem}
\usepackage{hyperref}

\geometry{left=2.5cm,right=2.5cm,top=2.4cm,bottom=2.4cm}
\setlist{nosep}
\hypersetup{colorlinks=true,linkcolor=black,urlcolor=blue}

\title{B 题\quad 高精度非圆曲线磨削的砂轮中心轨迹规划}
\author{}
\date{}

\begin{document}
\maketitle

\section*{背景}

非圆曲线零件广泛存在于凸轮、非圆轴承滚道、异形密封槽和精密泵传动机构中。此类零件的轮廓误差会直接影响传动稳定性、密封性能、噪声振动和寿命。对于高精度磨削，工件理论轮廓通常由离散坐标点给出，而数控机床实际执行的是砂轮中心轨迹。因此，需要把离散工件轮廓转换为可加工、可插补、可评价的砂轮中心刀位点序列。

本题考虑半径为 $R$ 的圆弧砂轮对平面非圆轮廓进行外接触磨削。理想情况下，砂轮中心应沿着与工件轮廓保持固定距离的轨迹运动。实际建模中必须处理离散点重构、法向和曲率估计、刀位点生成、等间距插补、机床动态约束、几何干涉和误差补偿等问题。

本题附件给出了某非圆凸轮轮廓的原始离散点和推荐建模点。原始数据可用于数据质量检查；除非特别说明，问题二至问题四均以推荐建模点为默认输入。

\section*{问题}

\subsection*{问题 1\quad 轮廓数据检查与高精度重构}

附件提供原始轮廓点和推荐建模点。请首先分析两组数据的闭合性、点序方向、坐标范围、相邻点间距和可能的异常点，说明推荐建模点相对原始点的合理性。然后建立工件轮廓 $C_w(t)=(x(t),y(t))$ 的光滑参数化模型。

要求回答：
\begin{enumerate}
  \item 采用何种参数化方式，如累积弦长参数、弧长参数、周期样条、NURBS 或三角插值；
  \item 如何保证封闭曲线的首尾连续性，至少说明 $C^1$ 或 $C^2$ 连续条件；
  \item 计算拟合误差、总弧长和反映局部弯曲程度的几何量；
  \item 讨论数据质量、采样密度和数值微分对微米级加工精度的影响。
\end{enumerate}

\subsection*{问题 2\quad 砂轮中心等距轨迹建模}

设砂轮半径 $R=10$ mm。请基于问题 1 得到的连续工件轮廓，建立砂轮中心轨迹模型，使砂轮与工件轮廓保持合理的外接触关系。

要求回答：
\begin{enumerate}
  \item 说明砂轮中心轨迹的几何构造方法，并讨论点列方向对轨迹方向判断的影响；
  \item 比较“直接对离散点做局部法向偏移”和“先连续重构再生成等距线”两种方法的误差机理；
  \item 输出 $N=2000$ 个砂轮中心刀位点，结果 CSV 至少包含刀位点编号和砂轮中心 $x,y$ 坐标三列，字段名见提交要求；
  \item 分析等距线在曲率较大区域的数值稳定性。
\end{enumerate}

\subsection*{问题 3\quad 等弧长刀位点规划与进给约束}

数控磨削希望砂轮中心刀位点沿轨迹近似等弧长分布，以获得更稳定的进给速度和表面质量。设砂轮中心轨迹总弧长为 $L_o$，理想相邻弧长为 $L_o/N$。

请建立等弧长离散化算法，使刀位点 $Q_0,Q_1,\ldots,Q_{N-1}$ 沿砂轮中心轨迹尽量均匀分布。
要求回答：
\begin{enumerate}
  \item 建立从曲线参数到弧长位置的数值转换方法；
  \item 给出 $N=1000,2000,5000,10000$ 时的弧长误差、计算耗时和精度收益；
  \item 在进给速度 $v_f=60$ mm/min、插补周期 $T_s=0.5$ ms、最大加速度和加加速度受限条件下，分析速度波动率是否满足 $0.1\%$ 目标；
  \item 若局部曲率变化过大，请提出局部加密、前瞻减速或平滑过渡策略。
\end{enumerate}

\subsection*{问题 4\quad 干涉检测、砂轮选型与误差补偿}

实际加工中，砂轮半径过大可能导致几何干涉；砂轮磨损和机床定位误差又会造成批量加工误差累积。请在问题 1--3 的基础上建立面向工程应用的最终方案。

要求回答：
\begin{enumerate}
  \item 对候选砂轮半径 $R\in\{5,8,10,15,20\}$ mm 建立干涉判定模型，定位可能过切的轮廓区段；
  \item 以轮廓度误差、材料去除效率和工艺可实现性为目标，选择合适砂轮半径或分区加工方案；
  \item 考虑机床重复定位误差 $\pm0.8\,\mu$m 和砂轮半径每件磨损 $2\,\mu$m，建立在线误差补偿模型；
  \item 综合给出最终加工轨迹规划流程图、公式或伪代码，并用 \texttt{evaluate\_submission.py} 检查提交刀位点的等距误差和等弧长均匀性。
\end{enumerate}

\section*{名词说明}

\begin{enumerate}
  \item \textbf{工件轮廓} $C_w$：待加工零件的理论边界曲线。附件中以离散坐标点形式给出，坐标单位为 mm。
  \item \textbf{砂轮中心轨迹}：圆弧砂轮中心在机床坐标系中的运动轨迹。对于半径为 $R$ 的砂轮，理想情况下它应与工件轮廓保持正确的几何接触关系。
  \item \textbf{等距曲线}：与原曲线保持固定距离的一类曲线。直接对离散点平移通常不等价于连续曲线意义下的等距轨迹。
  \item \textbf{刀位点}：数控系统实际执行的离散砂轮中心坐标点。提交结果中的每一行对应一个刀位点。
  \item \textbf{等弧长离散}：相邻刀位点沿砂轮中心轨迹的弧长尽量相等，用于减小进给速度波动。
  \item \textbf{干涉或过切}：当砂轮半径相对局部曲率半径过大时，砂轮会切入不应加工的区域，导致轮廓破坏。
  \item \textbf{轮廓度误差}：实际加工轮廓相对理论轮廓的最大几何偏差，本题目标量级为微米级。
\end{enumerate}

\section*{评价指标}

公开评价脚本主要报告两类代理指标：一类衡量刀位点到推荐建模轮廓的距离是否接近砂轮半径，另一类衡量相邻刀位点间距是否均匀。严格评价时应使用连续重构曲线和真实弧长，而非仅使用折线距离和弦长。公开脚本主要用于格式检查和结果自测，不能替代正式高精度评价。

问题目标为：轮廓重构误差小于 $0.1\,\mu$m，刀位点等弧长误差尽量小于 $0.1\,\mu$m，最终轮廓度误差尽量小于 $1\,\mu$m，并给出可解释的干涉和补偿方案。

\section*{数据说明}

本题提供以下附件：
\begin{enumerate}
  \item \texttt{attachment1\_profile\_points\_raw.csv}：原始轮廓离散点；
  \item \texttt{attachment2\_profile\_points\_clean.csv}：推荐建模轮廓点，问题二至问题四默认使用；
  \item \texttt{attachment3\_machine\_params.csv}：砂轮、机床和精度目标参数；
  \item \texttt{attachment4\_baseline\_summary.csv}：基础数据统计和样例提交评价结果；
  \item \texttt{attachment5\_data\_dictionary.csv}：附件字段说明；
  \item \texttt{attachment6\_problem\_diagram.png}：工件轮廓与砂轮中心轨迹关系示意图；
  \item \texttt{sample\_submission\_discrete\_offset\_2000.csv}：样例提交文件；
  \item \texttt{evaluate\_submission.py}：公开评价脚本。
\end{enumerate}

注意：样例提交仅用于说明格式，不代表最优方法或最终结果。推荐建模点坐标单位为 mm，计算过程中不宜过早四舍五入。

\section*{提交要求}

参赛队应提交建模论文、结果文件和可复现结果的支撑代码。结果文件为 CSV 格式，至少包含三列：
\[
\texttt{tool\_id},\qquad \texttt{x\_center\_mm},\qquad \texttt{y\_center\_mm}.
\]

论文中应包括：模型假设、变量说明、公式推导、参数估计、误差分析、干涉判定、补偿方案、数值结果和模型局限性分析。支撑代码要求参照全国大学生数学建模竞赛对支撑材料的要求：提交完整源程序、运行说明和依赖环境说明；代码应能从本题附件数据出发，复现论文中的主要图表、最终提交 CSV 和主要评价指标；路径尽量使用相对路径，不依赖本机私有目录、手工中间操作或未随题提供的数据。若算法含随机过程，应固定随机种子并说明可能差异。

\end{document}
